%  LaTeX support: latex@mdpi.com 
%  For support, please attach all files needed for compiling as well as the log file, and specify your operating system, LaTeX version, and LaTeX editor.

%=================================================================
\documentclass[preprints,article,submit,pdftex,oneauthor]{Definitions/mdpi} 
%\documentclass[preprints,article,submit,pdftex,moreauthors]{Definitions/mdpi} 
% For posting an early version of this manuscript as a preprint, you may use "preprints" as the journal. Changing "submit" to "accept" before posting will remove line numbers.

% Below journals will use APA reference format:
% admsci, aieduc, behavsci, businesses, econometrics, economies, education, ejihpe, famsci, games, humans, ijcs, ijfs, jintelligence, journalmedia, jrfm, jsam, languages, peacestud, psycholint, publications, tourismhosp, youth

% Below journals will use Chicago reference format:
% arts, genealogy, histories, humanities, laws, literature, religions, risks, socsci

%--------------------
% Class Options:
%--------------------
%----------
% journal
%----------
% Choose between the following MDPI journals:
% accountaudit, acoustics, actuators, addictions, adhesives, admsci, adolescents, aerobiology, aerospace, agriculture, agriengineering, agrochemicals, agronomy, ai, aichem, aieduc, aieng, aimater, aimed, aipa, air, aisens, algorithms, allergies, alloys, amh, analog, analytica, analytics, anatomia, anesthres, animals, antibiotics, antibodies, antioxidants, applbiosci, appliedchem, appliedmath, appliedphys, applmech, applmicrobiol, applnano, applsci, aquacj, architecture, arm, arthropoda, arts, asc, asi, astronautics, astronomy, atmosphere, atoms, audiolres, automation, axioms, bacteria, batteries, bdcc, behavsci, beverages, biochem, bioengineering, biologics, biology, biomass, biomechanics, biomed, biomedicines, biomedinformatics, biomimetics, biomolecules, biophysica, bioresourbioprod, biosensors, biosphere, biotech, birds, blockchains, bloods, blsf, brainsci, breath, buildings, businesses, cancers, carbon, cardio, cardiogenetics, cardiovascmed, catalysts, cells, ceramics, challenges, chemengineering, chemistry, chemosensors, chemproc, children, chips, cimb, civileng, cleantechnol, climate, clinbioenerg, clinpract, clockssleep, cmd, cmtr, coasts, coatings, colloids, colorants, commodities, complexities, complications, compounds, computation, computers, condensedmatter, conservation, constrmater, cosmetics, covid, crops, cryo, cryptography, crystals, csmf, ctn, culture, curroncol, cyber, dairy, data, ddc, dentistry, dermato, dermatopathology, designs, devices, dhi, diabetology, diagnostics, dietetics, digital, disabilities, diseases, diversity, dna, drones, dynamics, earth, ebj, ecm, ecologies, econometrics, economies, edm, education, eesp, ejihpe, electricity, electrochem, electronicmat, electronics, encyclopedia, endocrines, energies, eng, engproc, entomology, entropy, environments, environremediat, epidemiologia, epigenomes, esa, est, famsci, fermentation, fibers, fintech, fire, fishes, fluids, foods, forecasting, forensicsci, forests, fossstud, foundations, fractalfract, fuels, future, futureinternet, futurepharmacol, futurephys, futuretransp, galaxies, games, gases, gastroent, gastrointestdisord, gastronomy, gels, genealogy, genes, geographies, geohazards, geomatics, geometry, geosciences, geotechnics, geriatrics, germs, glacies, grasses, green, greenhealth, gucdd, hardware, hazardousmatters, healthcare, hearts, hemato, hematolrep, hep, heritage, higheredu, highthroughput, histories, horticulturae, hospitals, humanities, humans, hydrobiology, hydrogen, hydrology, hydropower, hygiene, idr, iic, ijcs, ijem, ijerph, ijfs, ijgi, ijmd, ijms, ijns, ijom, ijpb, ijt, ijtm, ijtpp, ime, immuno, informatics, information, infrastructures, inorganics, insects, instruments, inventions, iot, j, jaestheticmed, jal, jcdd, jcm, jcp, jcrm, jcs, jcto, jdad, jdb, jdream, jemr, jeta, jfb, jfmk, jgbg, jgg, jimaging, jintelligence, jlpea, jmahp, jmmp, jmms, jmp, jmse, jne, jnt, jof, joi, joitmc, joma, jop, joptm, jor, journalmedia, jox, jpbi, jphytomed, jpm, jrfm, jsam, jsan, jtaer, jvd, jzbg, kidneydial, kinasesphosphatases, knowledge, labmed, laboratories, lae, land, languages, laws, life, lights, limnolrev, lipidology, liquids, literature, livers, logics, logistics, lubricants, lymphatics, machines, macromol, magnetism, magnetochemistry, make, marinedrugs, materials, materproc, mathematics, mca, measurements, medicina, medicines, medsci, membranes, merits, metabolites, metals, meteorology, methane, metrics, metrology, micro, microarrays, microbiolres, microelectronics, micromachines, microorganisms, microplastics, microwave, minerals, mining, mmphys, modelling, molbank, molecules, mps, msf, mti, multimedia, muscles, nanoenergyadv, nanomanufacturing, nanomaterials, ncrna, ndt, network, neuroglia, neuroimaging, neurolint, neurosci, nitrogen, notspecified, nursrep, nutraceuticals, nutrients, obesities, occuphealth, oceans, ohbm, onco, optics, oral, organics, organoids, osteology, oxygen, pandemics, parasites, parasitologia, particles, pathogens, pathophysiology, peacestud, pediatrrep, pets, pharmaceuticals, pharmaceutics, pharmacoepidemiology, pharmacy, philosophies, photochem, photonics, phycology, physchem, physics, physiologia, plants, plasma, platforms, pollutants, polymers, polysaccharides, populations, poultry, powders, precipitation, precisoncol, preprints, proceedings, processes, prosthesis, proteomes, psf, psychiatryint, psychoactives, psycholint, publications, purification, quantumrep, quaternary, qubs, radiation, rdt, reactions, realestate, receptors, recycling, regeneration, religions, remotesensing, reports, reprodmed, resources, rheumato, risks, rjpm, robotics, rsee, ruminants, safety, sci, scipharm, sclerosis, seeds, sensors, separations, sexes, shi, signals, sinusitis, siuj, skins, smartcities, sna, societies, socsci, software, soilsystems, solar, solids, spectroscj, sports, standards, stats, std, stratsediment, stresses, surfaces, surgeries, suschem, sustainability, symmetry, synbio, systems, tae, targets, taxonomy, technologies, telecom, test, textiles, thalassrep, therapeutics, thermo, timespace, tomography, tourismhosp, toxics, toxins, tph, transplantology, transportation, traumacare, traumas, tri, tropicalmed, universe, urbansci, uro, vaccines, vehicles, venereology, vetsci, vibration, virtualworlds, viruses, vision, waste, water, welding, wem, wevj, wild, wind, women, world, youth, zoonoticdis

%---------
% article
%---------
% The default type of manuscript is "article", but can be replaced by: 
% abstract, addendum, article, benchmark, book, bookreview, briefcommunication, briefreport, casereport, changes, clinicopathologicalchallenge, comment, commentary, communication, conceptpaper, conferenceproceedings, correction, conferencereport, creative, datadescriptor, discussion, entry, expressionofconcern, extendedabstract, editorial, essay, erratum, fieldguide, hypothesis, interestingimages, letter, meetingreport, monograph, newbookreceived, obituary, opinion, proceedingpaper, projectreport, reply, retraction, review, perspective, protocol, shortnote, studyprotocol, supfile, systematicreview, technicalnote, viewpoint, guidelines, registeredreport, tutorial,  giantsinurology, urologyaroundtheworld
% supfile = supplementary materials

%----------
% submit
%----------
% The class option "submit" will be changed to "accept" by the Editorial Office when the paper is accepted. This will only make changes to the frontpage (e.g., the logo of the journal will get visible), the headings, and the copyright information. Also, line numbering will be removed. Journal info and pagination for accepted papers will also be assigned by the Editorial Office.

%------------------
% moreauthors
%------------------
% If there is only one author the class option oneauthor should be used. Otherwise use the class option moreauthors.

%---------
% pdftex
%---------
% The option pdftex is for use with pdfLaTeX. Remove "pdftex" for (1) compiling with LaTeX & dvi2pdf (if eps figures are used) or for (2) compiling with XeLaTeX.

%=================================================================
% MDPI internal commands - do not modify
\firstpage{1} 
\makeatletter 
\setcounter{page}{\@firstpage} 
\makeatother
\pubvolume{1}
\issuenum{1}
\articlenumber{0}
\pubyear{2026}
\copyrightyear{2026}
%\externaleditor{Firstname Lastname} % More than 1 editor, please add `` and '' before the last editor name
\datereceived{ } 
\daterevised{ } % Comment out if no revised date
\dateaccepted{ } 
\datepublished{ } 
%\datecorrected{} % For corrected papers: "Corrected: XXX" date in the original paper.
%\dateretracted{} % For retracted papers: "Retracted: XXX" date in the original paper.
%\doinum{} % Used for some special journals, like molbank
%\pdfoutput=1 % Uncommented for upload to arXiv.org
%\CorrStatement{yes}  % For updates
%\longauthorlist{yes} % For many authors that exceed the left citation part
%\IsAssociation{yes} % For association journals

%=================================================================
% Add packages and commands here. The following packages are loaded in our class file: fontenc, inputenc, calc, indentfirst, fancyhdr, graphicx, epstopdf, lastpage, ifthen, float, amsmath, amssymb, lineno, setspace, enumitem, mathpazo, booktabs, titlesec, etoolbox, tabto, xcolor, colortbl, soul, multirow, microtype, tikz, totcount, changepage, attrib, upgreek, array, tabularx, pbox, ragged2e, tocloft, marginnote, marginfix, enotez, amsthm, natbib, hyperref, cleveref, scrextend, url, geometry, newfloat, caption, draftwatermark, seqsplit
% cleveref: load \crefname definitions after \begin{document}

%=================================================================
% Please use the following mathematics environments: Theorem, Lemma, Corollary, Proposition, Characterization, Property, Problem, Example, ExamplesandDefinitions, Hypothesis, Remark, Definition, Notation, Assumption
%% For proofs, please use the proof environment (the amsthm package is loaded by the MDPI class).

%=================================================================
% Full title of the paper (Capitalized)
\Title{The Marshall Triangle: A Triadic Geometric Framework for Modeling Balance in Complex Systems}

% Author Orchid ID: enter ID or remove command
\newcommand{\orcidauthorA}{0009-0004-3674-7883} % Add \orcidA{} behind the author's name
%\newcommand{\orcidauthorB}{0000-0000-0000-000X} % Add \orcidB{} behind the author's name

% Authors, for the paper (add full first names)
\Author{Paul W. Marshall $^{1,*}$\orcidA{}}

\AuthorNames{Paul W. Marshall}

\address{%
$^{1}$ \quad Fidelitas LLC -- Series 1; pmarshall@fidelitas.llc}

\corres{Correspondence: pmarshall@fidelitas.llc}

% Current address and/or shared authorship
%\firstnote{Current address: Affiliation.}  
% Current address should not be the same as any items in the Affiliation section.

%\secondnote{These authors contributed equally to this work.}
% The commands \thirdnote{} till \eighthnote{} are available for further notes.

%\simplesumm{} % Simple summary

%\conference{} % An extended version of a conference paper

% Abstract (Do not insert blank lines, i.e. \\) 
\abstract{This paper introduces the Marshall Triangle, a formal geometric framework for representing balance among three competing system priorities. Inspired by color-space visualization, simplex geometry, and multi-objective tradeoff modeling, the framework maps triadic states into a continuous triangular domain in which relative dominance, pairwise blending, and equilibrium can be explored visually and mathematically. The model places primary influences at edge midpoints rather than vertices, producing a field in which secondary states emerge at the vertices and a white point marks neutral recombination within the rendered domain. A reference Harmony Index is defined from the normalized distance between a triadic state and balanced equilibrium, while the visualization layer is expressed as a Gaussian radial field over the triangular domain. The framework is implementation-independent but is illustrated through a reference rendering model using additive RGB contributions from midpoint sources. The Marshall Triangle is intended as a conceptual and computational tool for system design, visualization, and exploratory analysis in domains where competing priorities must be compared, calibrated, and communicated.}

% Keywords
\keyword{triadic modeling; computational geometry; simplex visualization; system balance; harmony index; visual analytics; multi-objective tradeoffs; color-space models; radial basis functions}

% The fields PACS, MSC, and JEL may be left empty or commented out if not applicable
%\PACS{J0101}
%\MSC{}
%\JEL{}

%%%%%%%%%%%%%%%%%%%%%%%%%%%%%%%%%%%%%%%%%%
% Only for the journal Diversity
%\LSID{\url{http://}}

%%%%%%%%%%%%%%%%%%%%%%%%%%%%%%%%%%%%%%%%%%
% Only for the journal Applied Sciences
%\featuredapplication{Authors are encouraged to provide a concise description of the specific application or a potential application of the work. This section is not mandatory.}
%%%%%%%%%%%%%%%%%%%%%%%%%%%%%%%%%%%%%%%%%%

%%%%%%%%%%%%%%%%%%%%%%%%%%%%%%%%%%%%%%%%%%
% Only for the journal Data
%\dataset{DOI number or link to the deposited data set if the data set is published separately. If the data set shall be published as a supplement to this paper, this field will be filled by the journal editors. In this case, please submit the data set as a supplement.}
%\datasetlicense{License under which the data set is made available (CC0, CC-BY, CC-BY-SA, CC-BY-NC, etc.)}

%%%%%%%%%%%%%%%%%%%%%%%%%%%%%%%%%%%%%%%%%%
% Only for the journal BioTech, Fishes, Neuroimaging and Toxins
%\keycontribution{The breakthroughs or highlights of the manuscript. Authors can write one or two sentences to describe the most important part of the paper.}

%%%%%%%%%%%%%%%%%%%%%%%%%%%%%%%%%%%%%%%%%%
% Only for the journal Encyclopedia
%\encyclopediadef{For entry manuscripts only: please provide a brief overview of the entry title instead of an abstract.}

%%%%%%%%%%%%%%%%%%%%%%%%%%%%%%%%%%%%%%%%%%
% Different journals have different requirements. Please check the specific journal guidelines in the "Instructions for Authors" on the journal's official website.
%\addhighlights{yes}
%\renewcommand{\addhighlights}{%
%
%\noindent The goal is to increase the discoverability and readability of the article via search engines and other scholars. Highlights should not be a copy of the abstract, but a simple text allowing the reader to quickly and simplified find out what the article is about and what can be cited from it. Each of these parts should be devoted up to 2~bullet points.\vspace{3pt}\\
%\textbf{What are the main findings?}
% \begin{itemize}[labelsep=2.5mm,topsep=-3pt]
% \item First bullet.
% \item Second bullet.
% \end{itemize}\vspace{3pt}
%\textbf{What are the implications of the main findings?}
% \begin{itemize}[labelsep=2.5mm,topsep=-3pt]
% \item First bullet.
% \item Second bullet.
% \end{itemize}
%}

%%%%%%%%%%%%%%%%%%%%%%%%%%%%%%%%%%%%%%%%%%
\begin{document}
\iffalse
\section{Introduction}
%%%%%%%%%%%%%%%%%%%%%%%%%%%%%%%%%%%%%%%%%%
\setcounter{section}{-1} %% Remove this when starting to work on the template.
\section{How to Use this Template}

The template details the sections that can be used in a manuscript. Note that the order and names of article sections may differ from the requirements of the journal (e.g., the positioning of the Materials and Methods section). Please check the instructions on the authors' page of the journal to verify the correct order and names. For any questions, please contact the editorial office of the journal or support@mdpi.com. For LaTeX-related questions please contact latex@mdpi.com.%\endnote{This is an endnote.} % To use endnotes, please un-comment \printendnotes below (before References). Only journal Laws uses \footnote.

% The order of the section titles is different for some journals. Please refer to the "Instructions for Authors” on the journal homepage.

\section{Introduction}

The introduction should briefly place the study in a broad context and highlight why it is important. It should define the purpose of the work and its significance. The current state of the research field should be reviewed carefully and key publications cited. Please highlight controversial and diverging hypotheses when necessary. Finally, briefly mention the main aim of the work and highlight the principal conclusions. As far as possible, please keep the introduction comprehensible to scientists outside your particular field of research. Citing a journal paper \citep{ref-journal}.  Now citing a book reference \citep{ref-book1,ref-book2} or other reference types \citep{ref-unpublish,ref-url}. Please use the command \citep{ref-proceeding,ref-thesis} for the following MDPI journals, which use author--date citation: Administrative Sciences, Arts, Behavioral Sciences, Businesses, Econometrics, Economies, Education Sciences, European Journal of Investigation in Health, Psychology and Education, Games, Genealogy, Histories, Humanities, Humans, IJFS, Journal of Intelligence, Journalism and Media, JRFM, Languages, Laws, Literature, Psychology International, Publications, Religions, Risks, Social Sciences, Tourism and Hospitality, Youth. 

%%%%%%%%%%%%%%%%%%%%%%%%%%%%%%%%%%%%%%%%%%
\section{Materials and Methods}

Materials and Methods should be described with sufficient details to allow others to replicate and build on published results. Please note that publication of your manuscript implicates that you must make all materials, data, computer code, and protocols associated with the publication available to readers. Please disclose at the submission stage any restrictions on the availability of materials or information. New methods and protocols should be described in detail while well-established methods can be briefly described and appropriately cited.

Research manuscripts reporting large datasets that are deposited in a publicly avail-able database should specify where the data have been deposited and provide the relevant accession numbers. If the accession numbers have not yet been obtained at the time of submission, please state that they will be provided during review. They must be provided prior to publication.

Interventionary studies involving animals or humans, and other studies require ethical approval must list the authority that provided approval and the corresponding ethical approval code.

In this section, where applicable, authors are required to disclose details of how gen-erative artificial intelligence (GenAI) has been used in this paper (e.g., to generate text, data, or graphics, or to assist in study design, data collection, analysis, or interpretation). The use of GenAI for superficial text editing (e.g., grammar, spelling, punctuation, and formatting) does not need to be declared.

\begin{quote}
This is an example of a quote.
\end{quote}

%%%%%%%%%%%%%%%%%%%%%%%%%%%%%%%%%%%%%%%%%%
\section{Results}

This section may be divided by subheadings. It should provide a concise and precise description of the experimental results, their interpretation as well as the experimental conclusions that can be drawn.
\subsection{Subsection}
\subsubsection{Subsubsection}

Bulleted lists look like this:
\begin{itemize}
\item	First bullet;
\item	Second bullet;
\item	Third bullet.
\end{itemize}

Numbered lists can be added as follows:
\begin{enumerate}
\item	First item; 
\item	Second item;
\item	Third item.
\end{enumerate}

The text continues here.

\subsection{Figures, Tables and Schemes}

All figures and tables should be cited in the main text as Figure~\ref{fig1}, Table~\ref{tab1}, etc.

\begin{figure}[H]
%\isPreprints{\centering}{} % Only used for preprints
\includegraphics[width=4.0 cm]{Definitions/logo-mdpi}
\caption{This is a figure. Schemes follow the same formatting.\label{fig1}}
\end{figure}   
\unskip

\begin{table}[H] 
%\small % Change table font size
\caption{This is a table caption. Tables should be placed in the main text near to the first time they are~cited.\label{tab1}}
%\isPreprints{\centering}{} % Only used for preprints
\begin{tabularx}{\textwidth}{CCC}
\toprule
\textbf{Title 1}	& \textbf{Title 2}	& \textbf{Title 3}\\
\midrule
Entry 1		& Data			& Data\\
Entry 2		& Data			& Data \textsuperscript{1}\\
\bottomrule
\end{tabularx}

\noindent{\footnotesize{\textsuperscript{1} Tables may have a footer.}}
\end{table}

The text continues here (Figure~\ref{fig2} and Table~\ref{tab2}).

% Example of a figure that spans the whole page width and with subfigures. The same concept works for tables, too.
\begin{figure}[H]
%\isPreprints{}{% This command is only used for ``preprints''.
\begin{adjustwidth}{-\extralength}{0cm}
\centering
%} % If the paper is ``preprints'', please uncomment this parenthesis.
\subfloat[\centering]{\includegraphics[width=7.0cm]{Definitions/logo-mdpi}}
%\hfill
\subfloat[\centering]{\includegraphics[width=7.0cm]{Definitions/logo-mdpi}}\\
\subfloat[\centering]{\includegraphics[width=7.0cm]{Definitions/logo-mdpi}}
%\hfill
\subfloat[\centering]{\includegraphics[width=7.0cm]{Definitions/logo-mdpi}}
%\isPreprints{}{% This command is only used for ``preprints''.
\end{adjustwidth}
%} % If the paper is ``preprints'', please uncomment this parenthesis.
\caption{This is a wide figure. Schemes follow the same formatting. If there are multiple panels, they should be listed as: (\textbf{a}) Description of what is contained in the first panel. (\textbf{b}) Description of what is contained in the second panel. (\textbf{c}) Description of what is contained in the third panel. (\textbf{d}) Description of what is contained in the fourth panel. Figures should be placed in the main text near to the first time they are cited. A caption on a single line should be centered.\label{fig2}}
\end{figure} 

\begin{table}[H]
\caption{This is a wide table.\label{tab2}}
%\isPreprints{\centering}{% This command is only used for ``preprints''.
	\begin{adjustwidth}{-\extralength}{0cm}
%} % If the paper is ``preprints'', please uncomment this parenthesis.
%\isPreprints{\begin{tabularx}{\textwidth}{CCCC}}{% This command is only used for ``preprints''.
		\begin{tabularx}{\fulllength}{CCCC}
%} % If the paper is ``preprints'', please uncomment this parenthesis.
			\toprule
			\textbf{Title 1}	& \textbf{Title 2}	& \textbf{Title 3}     & \textbf{Title 4}\\
			\midrule
\multirow[m]{3}{*}{Entry 1 *}	& Data			& Data			& Data\\
			  	                   & Data			& Data			& Data\\
			             	      & Data			& Data			& Data\\
                   \midrule
\multirow[m]{3}{*}{Entry 2}    & Data			& Data			& Data\\
			  	                  & Data			& Data			& Data\\
			             	     & Data			& Data			& Data\\
			\bottomrule
		\end{tabularx}
%		\isPreprints{}{% This command is only used for ``preprints''.
	\end{adjustwidth}
%} % If the paper is ``preprints'', please uncomment this parenthesis.
	\noindent{\footnotesize{* Tables may have a footer.}}
\end{table}

%\begin{listing}[H]
%\caption{Title of the listing}
%\rule{\columnwidth}{1pt}
%\raggedright Text of the listing. In font size footnotesize, small, or normalsize. Preferred format: left aligned and single spaced. Preferred border format: top border line and bottom border line.
%\rule{\columnwidth}{1pt}
%\end{listing}

Text.

Text.

\subsection{Formatting of Mathematical Components}

This is the example 1 of equation:
\begin{linenomath}
\begin{equation}
a = 1,
\end{equation}
\end{linenomath}
the text following an equation need not be a new paragraph. Please punctuate equations as regular text.
%% If the documentclass option "submit" is chosen, please insert a blank line before and after any math environment (equation and eqnarray environments). This ensures correct linenumbering. The blank line should be removed when the documentclass option is changed to "accept" because the text following an equation should not be a new paragraph.

This is the example 2 of equation:
%\isPreprints{}{% This command is only used for ``preprints''.
\begin{adjustwidth}{-\extralength}{0cm}
%} % If the paper is ``preprints'', please uncomment this parenthesis.
\begin{equation}
a = b + c + d + e + f + g + h + i + j + k + l + m + n + o + p + q + r + s + t + u + v + w + x + y + z
\end{equation}
%\isPreprints{}{% This command is only used for ``preprints''.
\end{adjustwidth}
%} % If the paper is ``preprints'', please uncomment this parenthesis.

%% Example of a page in landscape format (with table and table footnote).
%\startlandscape
%\begin{table}[H] %% Table in wide page
%%\isPreprints{\centering}{} % This command is only used for ``preprints''.
%\caption{This is a very wide table.\label{tab3}}
%	\begin{tabularx}{\textwidth}{CCCC}
%		\toprule
%		\textbf{Title 1}	& \textbf{Title 2}	& \textbf{Title 3}	& \textbf{Title 4}\\
%		\midrule
%		Entry 1		& Data			& Data			& This cell has some longer content that runs over two lines.\\
%		Entry 2		& Data			& Data			& Data\textsuperscript{1}\\
%		\bottomrule
%	\end{tabularx}
%%\isPreprints{}{% This command is only used for ``preprints''.
%	\begin{adjustwidth}{+\extralength}{0cm}
%%} % If the paper is ``preprints'', please uncomment this parenthesis.
%		\noindent\footnotesize{\textsuperscript{1} This is a table footnote.}
%%\isPreprints{}{% This command is only used for ``preprints''.
%	\end{adjustwidth}
%%} % If the paper is ``preprints'', please uncomment this parenthesis.
%\end{table}
%\finishlandscape

Please punctuate equations as regular text. Theorem-type environments (including propositions, lemmas, corollaries etc.) can be formatted as follows:
%% Example of a theorem:
\begin{Theorem}
Example text of a theorem.
\end{Theorem}

The text continues here. Proofs must be formatted as follows:

%% Example of a proof:
\begin{proof}[Proof of Theorem 1]
Text of the proof. Note that the phrase ``of Theorem 1'' is optional if it is clear which theorem is being referred to.
\end{proof}
The text continues here.

%%%%%%%%%%%%%%%%%%%%%%%%%%%%%%%%%%%%%%%%%%
\section{Discussion}

Authors should discuss the results and how they can be interpreted from the perspective of previous studies and of the working hypotheses. The findings and their implications should be discussed in the broadest context possible. Future research directions may also be highlighted.

%%%%%%%%%%%%%%%%%%%%%%%%%%%%%%%%%%%%%%%%%%
\section{Conclusions}

This section is not mandatory, but can be added to the manuscript if the discussion is unusually long or complex.

%%%%%%%%%%%%%%%%%%%%%%%%%%%%%%%%%%%%%%%%%%
\section{Patents}

This section is not mandatory, but may be added if there are patents resulting from the work reported in this manuscript.

%%%%%%%%%%%%%%%%%%%%%%%%%%%%%%%%%%%%%%%%%%
\vspace{6pt} 

%%%%%%%%%%%%%%%%%%%%%%%%%%%%%%%%%%%%%%%%%%
%% optional
%\supplementary{The following supporting information can be downloaded at:  \linksupplementary{s1}, Figure S1: title; Table S1: title; Video S1: title.}

% Only for journal Methods and Protocols:
% If you wish to submit a video article, please do so with any other supplementary material.
% \supplementary{The following supporting information can be downloaded at: \linksupplementary{s1}, Figure S1: title; Table S1: title; Video S1: title. A supporting video article is available at doi: link.}

% Only used for preprtints:
% \supplementary{The following supporting information can be downloaded at the website of this paper posted on \href{https://www.preprints.org/}{Preprints.org}.}

% Only for journal Hardware:
% If you wish to submit a video article, please do so with any other supplementary material.
% \supplementary{The following supporting information can be downloaded at: \linksupplementary{s1}, Figure S1: title; Table S1: title; Video S1: title.\vspace{6pt}\\
%\begin{tabularx}{\textwidth}{lll}
%\toprule
%\textbf{Name} & \textbf{Type} & \textbf{Description} \\
%\midrule
%S1 & Python script (.py) & Script of python source code used in XX \\
%S2 & Text (.txt) & Script of modelling code used to make Figure X \\
%S3 & Text (.txt) & Raw data from experiment X \\
%S4 & Video (.mp4) & Video demonstrating the hardware in use \\
%... & ... & ... \\
%\bottomrule
%\end{tabularx}
%}

%%%%%%%%%%%%%%%%%%%%%%%%%%%%%%%%%%%%%%%%%%
\fi
\section{Introduction}

Complex systems frequently require the reconciliation of competing priorities rather than the optimization of a single objective. In fields ranging from software architecture to interactive system design, practitioners must navigate tradeoffs among factors that cannot be simultaneously maximized. Conventional scalar metrics often obscure these tradeoffs, while high-dimensional models can become difficult to interpret or communicate.

The Marshall Triangle is proposed as a compact geometric framework for expressing triadic balance. Drawing inspiration from color-space visualization, simplex-based coordinate systems, and multi-objective tradeoff modeling, the framework encodes three competing axes within a triangular domain. This representation supports intuitive visualization of dominance, hybridization, and equilibrium, while remaining amenable to mathematical description and computational implementation.

The paper makes three contributions. First, it defines a midpoint-source triangular construction in which primary influences are positioned at edge midpoints rather than vertices. Second, it introduces a reference Harmony Index for quantifying deviation from triadic balance. Third, it describes a Gaussian radial-field rendering model that produces continuous visual gradients and a field-derived white point of neutral recombination. The result is a general-purpose visualization framework for exploring competing priorities in complex systems.

\section{Related Work and Conceptual Context}

The Marshall Triangle is related to, but distinct from, several established traditions. Colorimetry and color-space modeling provide a precedent for representing perceptual relationships through geometric domains and for treating white or neutral points as meaningful locations within a color system \citep{maxwell1860,cie1931,smithguild1931}. Information visualization and visual analytics similarly emphasize the use of visual structures to support analytical reasoning, comparison, and interpretation \citep{card1999,thomas2005,ware2012}.

The triangular domain also connects the framework to barycentric and simplex-based coordinate systems. In barycentric geometry, a point in a triangle may be represented by normalized weights relative to reference points, and these weights naturally express mixture, proportion, and relative contribution \citep{mobius1827,farin1990}. The Marshall Triangle adopts this general logic of normalized triadic representation, but its visual construction differs from a classical vertex-primary simplex by placing primary sources at edge midpoints.

The framework is also conceptually adjacent to multi-objective optimization, where tradeoffs arise because objectives often cannot be improved simultaneously without degrading others \citep{miettinen1999,deb2001}. However, the Marshall Triangle is not presented as an optimization algorithm or decision rule. It is an interpretive geometric and visual model for representing relationships among competing priorities, making tradeoff structure easier to inspect, calibrate, and communicate.

Finally, the reference implementation uses radial Gaussian falloff, connecting the visualization layer to radial basis function methods, in which local influence is modeled as a distance-dependent field around source points \citep{hardy1971,fasshauer2007}. This relation is conceptual rather than prescriptive: other interpolation or field models may be used if they preserve the core triadic interpretation.

\section{Conceptual Overview}

\subsection{Triadic Balance}

The Marshall Triangle is defined by three fundamental axes, each representing a competing priority within a system. In the reference application discussed throughout this paper, these axes correspond to Privacy, Performance, and Personalization. However, the framework itself is agnostic to the specific semantics of the axes; any triad of competing factors may be substituted.

In contrast to classical color triangles, the Marshall Triangle does not place primary axes at the vertices. Instead, the dominant source influences are located at the midpoints of each edge. The vertices therefore represent secondary dominance states emerging from the equal contribution of two axes, while interior points represent continuously blended states influenced by all three axes.

\subsection{Emergent Equilibrium}

The geometric center of the triangle corresponds to a symmetric balance among all three axes. However, the perceptual white point---the visual representation of neutral recombination---may shift based on calibration and weighting. This white point results from the recombination of triadic influences, dynamically emerging as a field-derived point of maximum coherence, even when the geometric center does not coincide with the calibrated perceptual balance.

\section{Geometric Construction}

The Marshall Triangle is constructed as an equilateral triangle embedded in a two-dimensional plane. Its defining innovation lies in relocating the primary source influences from the vertices---typical of classical color-triangle models---to the midpoints of each edge. Each midpoint anchors one of the three fundamental axes: Privacy, Performance, or Personalization. This inversion of conventional triangle logic establishes a continuous field of influence across the interior, rather than partitioned zones of dominance.

Within this configuration, any point inside the triangle represents a weighted combination of the three axes, constrained by non-negativity and normalization. The vertices emerge not as sources, but as extrema---locations where the contribution from two adjacent sources dominates while the third is minimized. This geometry creates a bounded, interpretable domain in which perceptual gradients, tension dynamics, and emergent equilibrium states can be visualized and calibrated with mathematical rigor.

\begin{figure}[H]
\centering
\includegraphics[width=0.75\textwidth]{Fig1_marshall_triangle-v1-sovereign.png}
\caption{The Marshall Triangle in its canonical visual form, showing a continuous triadic field structured by midpoint-based primary influences and an emergent white point of neutral recombination.}
\label{fig:canonical-triangle}
\end{figure}

\section{Mathematical Formulation}

Let a triadic system state be represented by a vector $s=(s_1,s_2,s_3)$, where each component is non-negative and at least one component is positive. The normalized state vector is

\begin{linenomath}
\begin{equation}
p_i = \frac{s_i}{\sum_{j=1}^{3} s_j}, \qquad i \in \{1,2,3\}.
\end{equation}
\end{linenomath}

The balanced reference state is $u=(1/3,1/3,1/3)$. A reference Harmony Index may be defined as

\begin{linenomath}
\begin{equation}
H(s)=1-\frac{\lVert p-u\rVert_2}{\sqrt{2/3}},
\end{equation}
\end{linenomath}

where $\sqrt{2/3}$ is the maximum Euclidean distance from the balanced state to a simplex vertex. Under this definition, $H=1$ at perfect triadic balance and approaches $0$ at maximal single-axis dominance.

For visualization, let $\Omega$ denote the triangular domain and let $c_1,c_2,c_3 \in \Omega$ denote the three edge-midpoint source locations. Distances on $\Omega$ are Euclidean distances in the planar embedding of the equilateral triangle, with all coordinates expressed in the same normalized units. For a point $x \in \Omega$, the radial influence of source $i$ may be written as

\begin{linenomath}
\begin{equation}
G_i(x)=\exp\left(-\frac{\lVert x-c_i\rVert^2}{2\sigma^2}\right).
\end{equation}
\end{linenomath}

The parameter $\sigma>0$ controls the radial falloff width and may be chosen as a calibrated fraction of the normalized triangle scale. Let $e_1,e_2,e_3$ denote the red, green, and blue basis channels. A rendered color field can then be expressed as

\begin{linenomath}
\begin{equation}
C(x;s)=N\left(\sum_{i=1}^{3} p_i G_i(x)e_i\right),
\end{equation}
\end{linenomath}

where $N(\cdot)$ is a hue-preserving normalization operator, for example $N(v)=v/\max(\lVert v\rVert_{\infty},\epsilon)$ for a small $\epsilon>0$, or an equivalent max-value normalization used to avoid RGB clipping. The reference implementation uses max-value normalization, but the framework does not require a single interpolation method. Valid implementations may employ barycentric coordinates, radial falloff fields, or other continuous blending strategies consistent with triadic equilibrium.

\section{Illustrative Examples}

Figure~\ref{fig:canonical-triangle} presents the Marshall Triangle in its canonical visual form. Figure~\ref{fig:labeled-structure} identifies the primary and pairwise regions of the triangle, distinguishing midpoint-based primary influences from vertex-adjacent blended regions. Figure~\ref{fig:state-comparison} then compares three representative system states: a dual-axis emphasis state, a balanced default state, and a single-axis Privacy-emphasis state.

\begin{figure}[H]
\centering
\includegraphics[width=0.90\textwidth]{Fig2_MT_Labels_Dark.png}
\caption{Labeled structural interpretation of the Marshall Triangle. Primary influences are positioned at edge midpoints, while vertex-adjacent regions express pairwise blending between adjacent axes. The interior field represents continuously blended triadic states, with the white point marking computed neutral recombination within the field.}
\label{fig:labeled-structure}
\end{figure}

Across these examples, the white point should not be read as a pointer toward the dominant axis. The white point is field-derived rather than axis-indicative: it marks the computed location of neutral recombination within the resulting color field. When one axis expands, the white point may shift away from that axis even as the corresponding field becomes more visually dominant across the triangle. When two axes expand together, the white point may shift toward the relatively minimized third axis. In both cases, the movement reflects the recalculated field geometry rather than a separate compensatory mechanism.

This construction also distinguishes primary-axis emphasis from pairwise blending. Near each edge midpoint, the corresponding primary axis is emphasized; near the vertices, adjacent axes blend into secondary dominance regions. The interior of the triangle represents continuously blended triadic states, where smooth gradients reflect varying relationships among Privacy, Performance, and Personalization.

\begin{figure}[H]
\centering
\includegraphics[width=0.92\textwidth]{Fig3_marshall_triangle_state_comparison.png}
\caption{Representative Marshall Triangle state comparison. (\textbf{a}) Privacy + Personalization emphasis, with Performance comparatively minimized. (\textbf{b}) Balanced default state. (\textbf{c}) Privacy-emphasis state. These panels illustrate that the white point is field-derived rather than axis-indicative: it marks the computed location of neutral recombination within the resulting color field, not the location of the dominant influence.}
\label{fig:state-comparison}
\end{figure}

These examples are intended to demonstrate the expressive behavior of the framework rather than prescribe a single computational implementation. Small adjustments to triadic weighting produce coherent visual transitions, enabling calibration, interpretation, and comparison of system states within a shared geometric field.

\section{Discussion}

The Marshall Triangle offers a middle ground between oversimplified scalar metrics and unwieldy high-dimensional models. By constraining analysis to a triadic domain, it foregrounds tradeoffs and balance in a visually and mathematically tractable form.

The framework is not intended to replace formal multi-objective optimization, empirical user modeling, or domain-specific decision analysis. Instead, it provides a compact representational layer for exploring and communicating triadic relationships before, during, or after more specialized analysis. Its interpretive value lies in making competing priorities visible as continuous fields rather than isolated scores.

The reference implementation adopts radial Gaussian falloff calibrated to an emergent white point. This implementation choice supports smooth visual transitions and intuitive manipulation, but the underlying framework remains broader than any single software realization. Future implementations may adapt the interpolation method, calibration procedure, or semantic axes while preserving the central midpoint-source construction.

\section{Limitations and Future Work}

The framework is intentionally formal and conceptual. It does not by itself establish empirical validity, normative priority weights, or domain-specific decision authority. The Privacy, Performance, and Personalization axes used in the reference application are illustrative rather than exhaustive.

Future work may include empirical studies of interpretability, formal comparison with barycentric and radial basis function interpolation methods, applications to additional triadic domains, and integration with decision-support systems. Extensions to higher-order simplices may also allow analogous visualization of four or more competing priorities, though such extensions would sacrifice some of the visual immediacy of the triangular form.

\section{Conclusions}

The Marshall Triangle provides a concise geometric framework for reasoning about balance among three competing priorities. By combining interpretability, mathematical structure, and visual expressiveness, it offers a flexible conceptual tool for analyzing complex systems. Its abstraction from specific implementations allows it to serve as a unifying model across domains where triadic tradeoffs are fundamental.

\authorcontributions{Conceptualization, software, visualization, formal analysis, and manuscript preparation were conducted solely by the author, Paul W. Marshall.}

\funding{This research received no external funding.}

\institutionalreview{Not applicable.}

\informedconsent{Not applicable.}

\dataavailability{The source code, figures, manuscript source files, and public release artifacts are available through the project repository at \url{https://github.com/Paul-W-Marshall/marshall-triangle}. The interactive implementation is available at \url{https://marshalltriangle.app}, and the explanatory landing site is available at \url{https://marshalltriangle.com}. An archival Zenodo DOI will be added following public release. The Story Protocol provenance record is available at \url{https://explorer.story.foundation/ipa/0x8eE1e758dafc5Cb0Ee6D862D26eAF68eE33cf920}.}

%\durcstatement{Current research is limited to the [please insert a specific academic field, e.g., XXX], which is beneficial [share benefits and/or primary use] and does not pose a threat to public health or national security. Authors acknowledge the dual-use potential of the research involving xxx and confirm that all necessary precautions have been taken to prevent potential misuse. As an ethical responsibility, authors strictly adhere to relevant national and international laws about DURC. Authors advocate for responsible deployment, ethical considerations, regulatory compliance, and transparent reporting to mitigate misuse risks and foster beneficial outcomes.}

% Only for journal Nursing Reports
%\publicinvolvement{Please describe how the public (patients, consumers, carers) were involved in the research. Consider reporting against the GRIPP2 (Guidance for Reporting Involvement of Patients and the Public) checklist. If the public were not involved in any aspect of the research add: ``No public involvement in any aspect of this research''.}
%
%% Only for journal Nursing Reports
%\guidelinesstandards{Please add a statement indicating which reporting guideline was used when drafting the report. For example, ``This manuscript was drafted against the XXX (the full name of reporting guidelines and citation) for XXX (type of research) research''. A complete list of reporting guidelines can be accessed via the equator network: \url{https://www.equator-network.org/}.}
%
%% Only for journal Nursing Reports
%\useofartificialintelligence{Please describe in detail any and all uses of artificial intelligence (AI) or AI-assisted tools used in the preparation of the manuscript. This may include, but is not limited to, language translation, language editing and grammar, or generating text. Alternatively, please state that “AI or AI-assisted tools were not used in drafting any aspect of this manuscript”.}

\acknowledgments{The author acknowledges the use of AI-assisted tools for editorial support, formatting assistance, and version tracking. The author reviewed and approved all content and takes full responsibility for the manuscript.}

\conflictsofinterest{The author is the originator of the Marshall Triangle framework and maintains related public project materials and software implementations. The author may have associated intellectual property interests in implementations of the framework. No external funding or third-party financial conflicts are declared.}

%%%%%%%%%%%%%%%%%%%%%%%%%%%%%%%%%%%%%%%%%%
%% Optional

%% Only for journal Encyclopedia
%\entrylink{The Link to this entry published on the encyclopedia platform.}


\reftitle{References}

\begin{thebibliography}{999}

\bibitem[Maxwell(1860)]{maxwell1860}
Maxwell, J.C. On the theory of compound colours, and the relations of the colours of the spectrum. \textit{Philosophical Transactions of the Royal Society of London} \textbf{1860}, \textit{150}, 57--84. https://doi.org/10.1098/rstl.1860.0005.

\bibitem[CIE(1932)]{cie1931}
Commission Internationale de l'Eclairage. \textit{Proceedings of the Eighth Session, Cambridge, 1931}; Cambridge University Press: Cambridge, UK, 1932.

\bibitem[Smith and Guild(1931)]{smithguild1931}
Smith, T.; Guild, J. The C.I.E. colorimetric standards and their use. \textit{Transactions of the Optical Society} \textbf{1931}, \textit{33}, 73--134.

\bibitem[Card et~al.(1999)]{card1999}
Card, S.K.; Mackinlay, J.D.; Shneiderman, B., Eds. \textit{Readings in Information Visualization: Using Vision to Think}; Morgan Kaufmann: San Francisco, CA, USA, 1999.

\bibitem[Thomas and Cook(2005)]{thomas2005}
Thomas, J.J.; Cook, K.A., Eds. \textit{Illuminating the Path: The Research and Development Agenda for Visual Analytics}; National Visualization and Analytics Center: Richland, WA, USA, 2005.

\bibitem[Ware(2012)]{ware2012}
Ware, C. \textit{Information Visualization: Perception for Design}, 3rd ed.; Morgan Kaufmann: Waltham, MA, USA, 2012.

\bibitem[M{\"o}bius(1827)]{mobius1827}
M{\"o}bius, A.F. \textit{Der barycentrische Calcul: Ein neues H{\"u}lfsmittel zur analytischen Behandlung der Geometrie}; Johann Ambrosius Barth: Leipzig, Germany, 1827.

\bibitem[Farin(1990)]{farin1990}
Farin, G. \textit{Curves and Surfaces for Computer-Aided Geometric Design: A Practical Guide}, 2nd ed.; Academic Press: Boston, MA, USA, 1990.

\bibitem[Miettinen(1999)]{miettinen1999}
Miettinen, K. \textit{Nonlinear Multiobjective Optimization}; Kluwer Academic Publishers: Boston, MA, USA, 1999.

\bibitem[Deb(2001)]{deb2001}
Deb, K. \textit{Multi-Objective Optimization Using Evolutionary Algorithms}; John Wiley \& Sons: Chichester, UK, 2001.

\bibitem[Hardy(1971)]{hardy1971}
Hardy, R.L. Multiquadric equations of topography and other irregular surfaces. \textit{Journal of Geophysical Research} \textbf{1971}, \textit{76}, 1905--1915.

\bibitem[Fasshauer(2007)]{fasshauer2007}
Fasshauer, G.E. \textit{Meshfree Approximation Methods with MATLAB}; World Scientific: Singapore, 2007.

\end{thebibliography}
% Please provide either the correct journal abbreviation (e.g. according to the “List of Title Word Abbreviations” http://www.issn.org/services/online-services/access-to-the-ltwa/) or the full name of the journal.
% Citations and References in Supplementary files are permitted provided that they also appear in the reference list here. 

%=====================================
% References, variant A: external bibliography
%=====================================
% \bibliography{your_external_BibTeX_file}

%=====================================
% References, variant B: internal bibliography
%=====================================

% ACS format%%%%%%%%%%%%
\enlargethispage{3\baselineskip}
\PublishersNote{}
\end{document}

